%% toc: P3 40 40.1
\chapter{Probability Spaces of Finite Structures}
\label{sec:p03-w03-p03-064-probability-spaces}

A random structure assigns, for every sufficiently large \(n\), a probability
distribution to the structures with exactly \(n\) elements.  This chapter
treats the questions that make such spaces a subject of logic: which
first-order properties hold with probability tending to one, at which
parameter values the limiting behaviour changes, which finite configurations
a sparse random graph forces to accompany a given set of vertices, and how a
game decides when the limiting theory is complete.  It works with the two
models of Spencer's survey, the simple unary predicate \(SU(n,p)\) and the
random graph \(G(n,p)\), and lets the parameter \(p=p(n)\) vary with \(n\).

\paragraph{Inputs and conclusions.}
The input is a model, \(SU(n,p(n))\) or \(G(n,p(n))\), a function \(p(n)\)
with values in \([0,1]\), and a first-order sentence \(A\) in the language
\(\{U\}\) or \(\{\sim\}\) with equality.  The conclusions are statements
about \(\lim_{n\to\infty}\Pr[\,\cdot\entails A]\): that it is \(0\) or \(1\),
where it changes as \(p(n)\) varies, or that it exists and lies strictly
between \(0\) and \(1\).  For the sparse graphs \(G(n,n^{-\alpha})\) with
\(\alpha\) irrational the chapter also gives the closure of a vertex set and
the almost sure bound on its size.  For theories without finite models it
gives a game criterion for completeness.  None of these conclusions decides
whether a particular finite structure has a property, and none decides a
property at a particular \(n\).

\paragraph{Source.}
The presentation follows Joel Spencer's \emph{Logic and Random Structures}
in the separately formatted author's version of 30 pages hosted on his New
York University page \cite{Spencer2007LogicRandomStructuresAuthor}.  Its
sections are numbered \(1\) to \(14\), followed by an unnumbered section
\emph{About the references}; every page locator below refers to this
version.  The printed chapter in \emph{Finite Model Theory and Its
Applications} (Springer, 2007) has not been consulted, and no locator here
refers to it.  Spencer's notation is kept: \(SU(n,p)\), \(G(n,p)\),
\(x\sim y\) for adjacency, \((R,H)\) for a rooted graph with parameters
\(r,v,e\), \(cl_t(X)\) for the \(t\)-closure and \(\mathrm{EHR}[G,H;t]\) for
the Ehrenfeucht game.  Where the source states a result without proof or with
an indication only, the annex proof is written for this chapter, and the
text says so.

\paragraph{Where the neighbouring material is.}
The random graph \(G(n,p(n))\), almost sure properties, the relations
\(\ll,\gg\) and threshold functions are defined in
\cref{def:p03-w03-p03-065-gnp,def:p03-w03-p03-065-almost-sure,def:p03-w03-p03-065-asymptotic-order,def:p03-w03-p03-065-threshold};
the same definitions apply verbatim to \(SU(n,p(n))\), whose definition is
\cref{def:p03-w03-p03-067-models}.  The clique, isolated-vertex and
common-neighbour thresholds, the Harris and Janson inequalities, rooted
graphs and the criterion that an extension statement holds almost surely
exactly when its rooted graph is safe are in
\cref{sec:p03-w03-p03-065-random-graphs}.  The zero--one law for
\(SU(n,p(n))\) when \(np(n)\to\infty\) and \(n(1-p(n))\to\infty\), with its
almost sure theory, completeness proof and countable model, is
\cref{thm:p03-w03-p03-067-unary}; the constant-\(p\) graph theory with the
Alice's Restaurant property and its countable model is
\cref{sec:p03-w03-p03-067-constant-p}.

\section{The unary model at its thresholds}
\label{sec:p03-w03-p03-064-unary}

Spencer opens with the sentence
\[
  \mathit{YES}:=\exists x\,U(x)
\]
in the language \(\{U\}\) with equality \cite[p.~1]{Spencer2007LogicRandomStructuresAuthor}.
It shows in the simplest case why \(p\) must be allowed to depend on \(n\).

\begin{hbtheorem}{thm:p03-w03-p03-064-yes}{proved=annex:p03-w03-p03-064-yes}{The sentence \(\mathit{YES}\) in \(SU(n,p)\)}
\begin{enumerate}
\item For every \(n\geq1\) and \(p\in[0,1]\),
  \(\Pr[SU(n,p)\entails\mathit{YES}]=1-(1-p)^n\); this is nondecreasing in
  \(p\), equal to \(0\) at \(p=0\) and to \(1\) at \(p=1\).
\item If \(p>0\) is constant, then
  \(\lim_{n\to\infty}\Pr[SU(n,p)\entails\mathit{YES}]=1\).
\item If \(c>0\) is constant and \(p(n)=c/n\) for \(n\geq c\), then
  \[
    \lim_{n\to\infty}\Pr[SU(n,p(n))\entails\mathit{YES}]=1-e^{-c}.
  \]
\item If \(p(n)\ll n^{-1}\), then \(\mathit{YES}\) holds almost never; if
  \(p(n)\gg n^{-1}\), then \(\mathit{YES}\) holds almost surely.  Thus
  \(n^{-1}\) is a threshold function for \(\mathit{YES}\).
\end{enumerate}
\end{hbtheorem}

These are Spencer's computations
\cite[pp.~1--2]{Spencer2007LogicRandomStructuresAuthor}.  For example, at
\(p(n)=n^{-1.01}\) the sentence holds almost never and at \(p(n)=n^{-0.99}\)
it holds almost surely \cite[p.~2]{Spencer2007LogicRandomStructuresAuthor}.
Item~3 shows the reparametrization \(p=c/n\) spreading the passage from
\(0\) to \(1\) over \(c\in(0,\infty)\); item~2 shows that the constant
parametrization misses it entirely.  The restriction \(n\geq c\) in item~3
only keeps \(p\leq1\), and limits ignore it.

Threshold functions are not unique (\cref{def:p03-w03-p03-065-threshold}):
\(5/n\) serves for \(\mathit{YES}\) as well as \(n^{-1}\).  Spencer adds that
a threshold function is determined roughly up to constant factors
\cite[p.~2]{Spencer2007LogicRandomStructuresAuthor}.  The following precise
form of that remark holds for every parametrized random structure.

\begin{hblemma}{thm:p03-w03-p03-064-threshold-ratio}{proved=annex:p03-w03-p03-064-threshold-ratio}{Two thresholds for one property}
Let \(p_0\) and \(p_1\) be threshold functions for the same property \(A\),
both positive for all large \(n\).  Then \(p_1/p_0\) does not tend to
\(\infty\) and does not tend to \(0\).
\end{hblemma}

The lemma does not say that \(p_1/p_0\) stays bounded; it excludes only
convergence to \(0\) or to \(\infty\).

Spencer's prototype zero--one law is stated for \(p(n)\gg n^{-1}\) and
\(1-p(n)\gg n^{-1}\) \cite[p.~3]{Spencer2007LogicRandomStructuresAuthor};
these are the hypotheses \(np(n)\to\infty\) and \(n(1-p(n))\to\infty\) of
\cref{thm:p03-w03-p03-067-unary}, which proves it through the theory of the
two schemas ``there are at least \(r\) points with \(U\)'' and ``there are at
least \(r\) points with \(\neg U\)''.  The hypotheses cannot be weakened to
include either threshold.

\begin{hbtheorem}{thm:p03-w03-p03-064-unary-boundary}{proved=annex:p03-w03-p03-064-unary-boundary}{The unary zero--one law fails at both thresholds}
Let \(c>0\) be constant.
\begin{enumerate}
\item If \(p(n)=c/n\), then
  \(\lim_n\Pr[SU(n,p(n))\entails\mathit{YES}]=1-e^{-c}\in(0,1)\).
\item If \(p(n)=1-c/n\), then
  \(\lim_n\Pr[SU(n,p(n))\entails\exists x\,\neg U(x)]=1-e^{-c}\in(0,1)\).
\end{enumerate}
In both regimes some first-order sentence has a limiting probability that is
neither \(0\) nor \(1\).
\end{hbtheorem}

\begin{hbwarning}{def:p03-w03-p03-064-limits-may-fail}{Three ways a threshold analysis can fail}
Spencer lists three difficulties that lead logicians to look for zero--one
laws rather than thresholds: a property need not be monotone in \(p\), a
threshold function need not exist, and the limit of the probabilities need
not exist \cite[pp.~2--3]{Spencer2007LogicRandomStructuresAuthor}.  The third
already occurs for \(\mathit{YES}\) if the parametrization is irregular: with
\(p(n)=1/n\) for even \(n\) and \(p(n)=2/n\) for odd \(n\), the proof of
\cref{thm:p03-w03-p03-064-yes} gives the two subsequential limits
\(1-e^{-1}\) and \(1-e^{-2}\).  A zero--one law asserts convergence as part of
its conclusion and holds only on a stated range of \(p(n)\).
\end{hbwarning}

\section{A threshold with a logarithmic factor: paths of length three}
\label{sec:p03-w03-p03-064-paths}

Spencer's last threshold example in \(G(n,p)\) is the sentence
\[
  \forall x_1\,\forall x_2\,\exists y_1\,\exists y_2\;
  (x_1\sim y_1\wedge y_1\sim y_2\wedge y_2\sim x_2),
\]
``every two vertices are joined by a path of length three'', for which he
gives the threshold function \(n^{-2/3}\ln^{1/3}n\)
\cite[p.~5]{Spencer2007LogicRandomStructuresAuthor}.  A fixed pair has
\(\binom{n-2}{2}\sim n^2/2\) unordered choices of the two inner vertices;
counting ordered choices, the expected number of paths is
\((n-2)(n-3)p^3\sim n^2p^3\).  That points to \(n^{-2/3}\), and the source
notes that a logarithmic factor is needed to reach \emph{every} pair.  The
displayed sentence has no clause \(x_1\neq x_2\).  Read literally at
\(x_1=x_2\) it would demand a triangle through every vertex, which is not the
property described in words; the guard \(x_1\neq x_2\) is therefore added
below.  The sentence also lets \(y_1\) or \(y_2\) coincide with \(x_1\) or
\(x_2\); for nonadjacent \(x_1,x_2\) this cannot happen, and the theorem is
stated for genuine paths so that both readings are covered.

\begin{hbdefinition}{def:p03-w03-p03-064-path3}{Paths of length three between two vertices}
In a graph, a \emph{path of length three} from \(x_1\) to \(x_2\) is a
sequence \(x_1,y_1,y_2,x_2\) of four pairwise distinct vertices with
\(x_1\sim y_1\), \(y_1\sim y_2\) and \(y_2\sim x_2\).  Let \(P_3\) be the
first-order property that every two distinct vertices are joined by a path
of length three, and \(P_3^{\mathrm w}\) the sentence
\(\forall x_1\,\forall x_2\,\bigl(x_1\neq x_2\rightarrow\exists y_1\,\exists y_2\,
(x_1\sim y_1\wedge y_1\sim y_2\wedge y_2\sim x_2)\bigr)\), in which
\(y_1,y_2\) need not be distinct from \(x_1,x_2\).
\end{hbdefinition}

\begin{hbtheorem}{thm:p03-w03-p03-064-path3}{proved=annex:p03-w03-p03-064-path3}{Threshold for paths of length three}
Let \(\varepsilon>0\) be fixed.
\begin{enumerate}
\item If \(n^2p^3\geq(2+\varepsilon)\ln n\) for all sufficiently large \(n\),
  then \(P_3\), and hence \(P_3^{\mathrm w}\), holds almost surely in
  \(G(n,p)\).
\item If \(\varepsilon<2\) and \(n^2p^3\leq(2-\varepsilon)\ln n\) for all
  sufficiently large \(n\), then almost surely some two nonadjacent vertices
  are joined by no path of length three; so \(P_3\) and
  \(P_3^{\mathrm w}\) hold almost never.
\end{enumerate}
Consequently \(n^{-2/3}\ln^{1/3}n\) is a threshold function for \(P_3\) and
for \(P_3^{\mathrm w}\).
\end{hbtheorem}

The proof applies the Janson inequality (\cref{thm:p03-w03-p03-065-janson})
to the paths of one pair, and to the paths of two pairs at once for a second
moment estimate.  The constant \(2\) in the theorem reflects the roughly
\(n^2/2\) pairs that must all succeed: a fixed pair fails with probability
about \(e^{-n^2p^3}\).

\section{Closure in sparse random graphs}
\label{sec:p03-w03-p03-064-closure}

Throughout this section \(\alpha\in(0,1)\) is irrational, and rooted graphs
\((R,H)\), their parameters \(v\) (nonroots) and \(e\) (edges not lying
between two roots), and the notions dense, sparse, rigid and safe are those
of \cref{def:p03-w03-p03-065-rooted-graph,def:p03-w03-p03-065-dense-sparse},
taken from Spencer's Section~3
\cite[p.~6]{Spencer2007LogicRandomStructuresAuthor}.  Rigid means that
\((S,H)\) is dense for every \(S\) with \(R\subseteq S\subsetneq V(H)\).
Graphs may be finite or infinite.

\begin{hbdefinition}{def:p03-w03-p03-064-rigid-extension}{Extensions inside a graph}
Let \(G\) be a graph, \(X\subseteq V(G)\), and \(Y\subseteq V(G)\) nonempty
and disjoint from \(X\), with \(X\) and \(Y\) finite.  Following Spencer,
vertices \(y_1,\ldots,y_v\) \emph{form an \((R,H)\) extension} over
\(x_1,\ldots,x_r\) if they have the adjacencies that \(H\) requires between
nonroots and between roots and nonroots, extra adjacencies being allowed; they
form a \emph{rigid extension} if they form an \((R,H)\) extension for some
rigid \((R,H)\) \cite[p.~8]{Spencer2007LogicRandomStructuresAuthor}.
Write \(G[X;Y]\) for the rooted graph \((X,G|_{X\cup Y})\), where
\(G|_{X\cup Y}\) is the induced subgraph.  For an arbitrary, possibly
infinite, \(X\subseteq V(G)\) and finite nonempty \(Y\) disjoint from \(X\),
\(Y\) is \emph{rigid over \(X\)} if, for every \(S\) with
\(X\subseteq S\subsetneq X\cup Y\),
\[
  |Y\setminus S|-e_G(S,Y)\,\alpha<0,
\]
where \(e_G(S,Y)\) is the number of edges of \(G\) with at least one end in
\(Y\setminus S\) and the other end in \(X\cup Y\).  For finite \(X\) this
says that \(G[X;Y]\) is rigid.
\end{hbdefinition}

\begin{hblemma}{thm:p03-w03-p03-064-rigid-basic}{proved=annex:p03-w03-p03-064-rigid-basic}{Rigid extensions: monotonicity and enlarging the base}
Let \(G\) be a graph.
\begin{enumerate}
\item If \((R,H)\) is rigid and \(H'\) is obtained from \(H\) by adding
  edges, then \((R,H')\) is rigid.  Consequently, for finite \(X\), a finite
  nonempty \(Y\) disjoint from \(X\) forms a rigid extension over \(X\)
  exactly when \(Y\) is rigid over \(X\).
\item If \(Y\) is rigid over \(X\), \(X\subseteq X'\subseteq V(G)\), and
  \(Y\not\subseteq X'\), then \(Y\setminus X'\) is rigid over \(X'\).
\item If \(Y\) is rigid over \(X\), then \(Y\) is rigid over every finite
  \(X_0\subseteq X\) that contains every vertex of \(X\) adjacent to a vertex
  of \(Y\).
\end{enumerate}
\end{hblemma}

\begin{hbdefinition}{def:p03-w03-p03-064-closure}{The \(t\)-closure}
Let \(t\geq1\), let \(G\) be a graph and \(X\subseteq V(G)\).  Call
\(Z\supseteq X\) \emph{\(t\)-closed} if no set of at most \(t\) vertices is
rigid over \(Z\).  The \(t\)-closure \(cl_t(X)\) is defined in two ways
\cite[p.~8]{Spencer2007LogicRandomStructuresAuthor}.
\begin{enumerate}
\item \emph{Algorithmically}: start with \(X\); whenever some \(Y\) with
  \(|Y|\leq t\) is rigid over the current set, add \(Y\) to it; \(cl_t(X)\)
  is the set of vertices that are added after finitely many such steps,
  together with \(X\).
\item \emph{Extremally}: \(cl_t(X)\) is the least \(t\)-closed set
  containing \(X\).
\end{enumerate}
For completeness the \(0\)-closure of \(X\) is \(X\) itself
\cite[p.~9]{Spencer2007LogicRandomStructuresAuthor}.  For tuples,
\(cl_t(x_1,\ldots,x_r)\cong cl_t(x_1',\ldots,x_r')\), for vertices of graphs
\(G\) and \(G'\), means that there is a graph isomorphism between the two
closures sending each \(x_i\) to \(x_i'\); and \(cl_t(x_1,\ldots,x_r)\cong H\)
means that the induced graph on the closure is isomorphic to \(H\) with the
roots \(x_1,\ldots,x_r\) in specified positions
\cite[pp.~8--9]{Spencer2007LogicRandomStructuresAuthor}.
\end{hbdefinition}

Spencer states that the two definitions agree, that the first does not
depend on the order in which rigid extensions are added, and that this rests
on combinatorial lemmas he omits
\cite[p.~8]{Spencer2007LogicRandomStructuresAuthor}.  The next theorem
supplies them.

\begin{hbtheorem}{thm:p03-w03-p03-064-closure-well-defined}{proved=annex:p03-w03-p03-064-closure-well-defined}{The two definitions of closure agree}
Let \(t\geq1\), \(G\) a graph and \(X\subseteq V(G)\).  Every sequence of
algorithmic steps from \(X\) stays inside every \(t\)-closed set containing
\(X\).  The set of all vertices reachable from \(X\) by finitely many steps,
together with \(X\), is \(t\)-closed.  Hence the least \(t\)-closed set
containing \(X\) exists, equals the algorithmic closure, and does not depend
on the order of the steps.  If \(G\) is finite, every maximal sequence of
steps is finite and ends exactly at \(cl_t(X)\).  Moreover
\(cl_t(X)\subseteq cl_{t'}(X)\) for \(t\leq t'\), and
\(cl_t(X)\subseteq cl_t(X')\) for \(X\subseteq X'\).
\end{hbtheorem}

\begin{hbexample}{def:p03-w03-p03-064-closure-example}{A \(4\)-closure at \(\alpha=\pi/7\)}
Spencer's example is that \(cl_4(x_1,x_2)\) might consist of \(x_1,x_2\);
\(y_1,y_2\), adjacent to each other and to both \(x_1\) and \(x_2\);
\(y_3,y_4,y_5,y_6\), forming a \(K_5\) with \(y_2\); and \(y_7\), a common
neighbour of \(x_2,y_1,y_5\) \cite[p.~8]{Spencer2007LogicRandomStructuresAuthor}.
With \(\alpha=\pi/7\approx0.4488\) each of the three steps is a rigid
extension with at most \(4\) nonroots.  For \(\{y_1,y_2\}\) over
\(\{x_1,x_2\}\): \(v=2\), \(e=5\), \(2-5\alpha<0\); fixing one of \(y_1,y_2\)
leaves \(v=1\), \(e=3\) and \(1-3\alpha<0\).  For \(\{y_3,\ldots,y_6\}\) over
the current set: the four new vertices span \(6\) edges and send \(4\) to
\(y_2\), so \(v=4\), \(e=10\), \(4-10\alpha<0\); fixing some of them leaves
\(v'\in\{1,2,3\}\) nonroots with \(e'=\binom{v'}2+v'(5-v')\in\{4,7,9\}\) edges,
and \(1-4\alpha\), \(2-7\alpha\), \(3-9\alpha\) are all negative.  For
\(\{y_7\}\): \(v=1\), \(e=3\), \(1-3\alpha<0\).  Extra edges only make these
numbers smaller (\cref{thm:p03-w03-p03-064-rigid-basic}).  So the nine
vertices lie in \(cl_4(x_1,x_2)\) whenever these adjacencies are present, and
by \cref{thm:p03-w03-p03-064-closure-well-defined} they are the whole
\(4\)-closure exactly when no set of at most four further vertices is rigid
over them.  The word ``might'' in the source refers to this last condition.
\end{hbexample}

\begin{hblemma}{thm:p03-w03-p03-064-closure-fo}{proved=annex:p03-w03-p03-064-closure-fo}{Closure types are first-order}
Fix \(t\geq0\), \(r\geq1\) and a finite graph \(H\) with distinguished
pairwise distinct vertices \(x_1,\ldots,x_r\).  There is a first-order
formula \(\theta_{t,H}(x_1,\ldots,x_r)\) in the language \(\{\sim\}\) such
that, in every graph \(G\) and for all pairwise distinct
\(a_1,\ldots,a_r\in V(G)\), \(G\entails\theta_{t,H}(\bar a)\) holds exactly
when \(cl_t(a_1,\ldots,a_r)\cong H\) with \(a_i\) in the position of
\(x_i\).
\end{hblemma}

Spencer records this observation and describes the formula for his example:
the existence of \(y_1,\ldots,y_7\) with the right adjacencies, together with
the nonexistence, for each of the finitely many rigid \((R,H)\) with
\(v\leq4\), of vertices extending \(x_1,\ldots,y_7\) in that way
\cite[p.~9]{Spencer2007LogicRandomStructuresAuthor}.  The proof in the annex
adds the condition that makes the formula exact: the specified vertices must
themselves be reachable by rigid steps.

\begin{hblemma}{thm:p03-w03-p03-064-nonexistence}{proved=annex:p03-w03-p03-064-nonexistence}{Nonexistence Lemma}
Let \(p=n^{-\alpha}\) and let \(M\geq1\).  Almost surely \(G(n,p)\) contains no
subgraph, induced or not, isomorphic to a graph with \(v\leq M\) vertices and
\(e\) edges such that \(v-e\alpha<0\).  In particular, for every \(t\geq1\),
almost surely \(cl_t(\emptyset)=\emptyset\) in \(G(n,n^{-\alpha})\)
\cite[p.~8]{Spencer2007LogicRandomStructuresAuthor}.
\end{hblemma}

\begin{hbtheorem}{thm:p03-w03-p03-064-finite-closure}{proved=annex:p03-w03-p03-064-finite-closure}{Finite Closure Lemma}
For all irrational \(\alpha\in(0,1)\) and integers \(r,t\geq1\) there is
\(K\) such that in \(G\sim G(n,n^{-\alpha})\) almost surely
\[
  |cl_t(x_1,\ldots,x_r)|<r+K\quad\text{for all }x_1,\ldots,x_r.
\]
One may take \(K=\lceil r/\epsilon\rceil\), where \(\epsilon\) is the least
value of \((e\alpha-v)/v\) over integers \(1\leq v\leq t\) and \(e\geq0\) with
\(v-e\alpha<0\) \cite[p.~9]{Spencer2007LogicRandomStructuresAuthor}.
\end{hbtheorem}

Spencer's proof treats the \(r\) roots as capital and each rigid step as a
cost of \(e\alpha-v\geq\epsilon v\).  A closure of size at least \(r+K\) would
contain a subgraph with fewer than \(r+K+t\) vertices whose own \(v-e\alpha\)
is negative, which the Nonexistence Lemma forbids
\cite[p.~9]{Spencer2007LogicRandomStructuresAuthor}.  Two details are
adjusted in the annex: a chain step may have \emph{at most} \(t\) nonroots,
and the final inequality is strict because \(\alpha\) is irrational, not
because \(r-K\epsilon\) is negative.  The minimum defining \(\epsilon\) exists
because, for each \(v\leq t\), only the least \(e\) with \(e\alpha>v\) matters.

\begin{hbwarning}{def:p03-w03-p03-064-closure-limits}{What the closure bounds do not say}
The bound \(K\) depends on \(\alpha\), \(r\) and \(t\), and it is an almost
sure statement about \(G(n,n^{-\alpha})\); a particular graph can have
arbitrarily large closures.  The closure \(cl(X)\), the union of the
\(cl_t(X)\) over all \(t\), is not first-order, and in a countable model of
the almost sure theory a single vertex can have infinite closure
(\cref{thm:p03-w03-p03-067-sparse-models}).  At rational \(\alpha\) the
quantity \(\epsilon\) can be \(0\), and the argument gives no bound.
\end{hbwarning}

\section{The Ehrenfeucht game and completeness}
\label{sec:p03-w03-p03-064-game}

Spencer's Section~14 introduces the game \(\mathrm{EHR}[G,H;t]\) for graphs
\(G,H\) and a positive integer \(t\): the board is a copy of each graph on
disjoint vertex sets; in each of \(t\) rounds Spoiler selects a vertex of
either graph and Duplicator a vertex of the other; \(x_i\in V(G)\) and
\(y_i\in V(H)\) are the vertices of round \(i\), whoever chose them; and
Duplicator wins exactly when \(x_i\sim x_j\iff y_i\sim y_j\) and
\(x_i=x_j\iff y_i=y_j\) for all \(i,j\)
\cite[p.~26]{Spencer2007LogicRandomStructuresAuthor}.  This is the game of
\cref{def:p03-ef-game} for the language \(\{\sim\}\)
(\cref{def:p03-w03-p03-067-game}).  For a finite relational language the
winning condition requires every relation symbol \(R\) of arity \(l\) to
have the same truth value on \((x_{i_1},\ldots,x_{i_l})\) and
\((y_{i_1},\ldots,y_{i_l})\) for every choice of \(i_1,\ldots,i_l\); this
covers the ordered graph and the unary language of \(SU(n,p)\)
\cite[p.~27]{Spencer2007LogicRandomStructuresAuthor}.  The basic theorem,
that Duplicator wins \(\mathrm{EHR}[G,H;t]\) if and only if \(G\) and \(H\)
satisfy the same first-order sentences of quantifier depth \(t\), is
\cref{thm:p03-ef} \cite[pp.~26--27]{Spencer2007LogicRandomStructuresAuthor}.

\begin{hblemma}{thm:p03-w03-p03-064-new-vertices}{proved=annex:p03-w03-p03-064-new-vertices}{Spoiler may be required to pick new vertices}
Let \(G,H\) be graphs with at least \(t\) vertices each.  Duplicator wins
\(\mathrm{EHR}[G,H;t]\) if and only if she wins the variant in which Spoiler
must always select a vertex not selected before in the graph he plays in.
In any winning play of either game, Duplicator answers a new vertex with a
new vertex.
\end{hblemma}

The remark is Spencer's \cite[p.~26]{Spencer2007LogicRandomStructuresAuthor}.
The hypothesis on the sizes is needed: if \(G\) has fewer than \(t\)
vertices, Spoiler in the variant may be unable to move.

\begin{hbexample}{def:p03-w03-p03-064-isolated-game}{An isolated vertex is seen in two rounds}
Let \(G\) have an isolated vertex and let \(H\) have none.  The sentence
\(\forall x\,\exists y\;x\sim y\) has quantifier depth \(2\) and separates
them, so by \cref{thm:p03-ef} Spoiler wins \(\mathrm{EHR}[G,H;2]\).
Spencer's direct strategy \cite[p.~27]{Spencer2007LogicRandomStructuresAuthor}:
Spoiler selects the isolated vertex \(x_1\in V(G)\), and Duplicator answers
some \(y_1\in V(H)\).  As \(y_1\) is not isolated, Spoiler selects in \(H\) a
vertex \(y_2\sim y_1\).  Duplicator must now answer \(x_2\in V(G)\) with
\(x_2\sim x_1\), and there is none: \(x_2=x_1\) fails because \(y_2\neq y_1\),
and every other \(x_2\) is nonadjacent to \(x_1\).
\end{hbexample}

\begin{hbcorollary}{thm:p03-w03-p03-064-ee-game}{proved=annex:p03-w03-p03-064-ee-game}{Elementary equivalence by finite games}
Two graphs \(G,H\) are elementarily equivalent if and only if Duplicator
wins \(\mathrm{EHR}[G,H;t]\) for every positive integer \(t\)
\cite[p.~27]{Spencer2007LogicRandomStructuresAuthor}.
\end{hbcorollary}

Spencer warns that this is not the same as Duplicator winning one game with
infinitely many moves \cite[p.~27]{Spencer2007LogicRandomStructuresAuthor}.
In the \emph{\(\omega\)-round game} the rounds are indexed by the positive
integers, and Duplicator wins if after every round the condition of
\(\mathrm{EHR}\) holds for the vertices chosen so far.

\begin{hbtheorem}{thm:p03-w03-p03-064-infinite-game}{proved=annex:p03-w03-p03-064-infinite-game}{Finite games and the infinite game differ}
\begin{enumerate}
\item If \(G\) and \(H\) are countable graphs and Duplicator has a winning
  strategy in the \(\omega\)-round game, then \(G\cong H\).
\item Let \(G\) be the disjoint union of complete graphs \(K_m\), one for
  each \(m\geq1\), and let \(H\) be the disjoint union of \(G\) and a
  countably infinite complete graph.  Then Duplicator wins
  \(\mathrm{EHR}[G,H;t]\) for every \(t\), so \(G\) and \(H\) are
  elementarily equivalent, but \(G\not\cong H\) and Duplicator has no
  winning strategy in the \(\omega\)-round game.
\end{enumerate}
\end{hbtheorem}

\begin{hbcorollary}{thm:p03-w03-p03-064-complete-game}{proved=annex:p03-w03-p03-064-complete-game}{Completeness by games on countable models}
Let \(T\) be a consistent first-order theory in the language \(\{\sim\}\),
or in any finite relational language, with no finite models.  Then \(T\) is
complete if and only if, for every two countable models \(G,H\) of \(T\) and
every positive integer \(t\), Duplicator wins \(\mathrm{EHR}[G,H;t]\)
\cite[p.~27]{Spencer2007LogicRandomStructuresAuthor}.
\end{hbcorollary}

This corollary is the completeness test Spencer uses for every almost sure
theory in the survey: for the unary theory of his Section~1
\cite[p.~3]{Spencer2007LogicRandomStructuresAuthor} and for the
constant-\(p\) graph theory of his Section~6
\cite[pp.~12--13]{Spencer2007LogicRandomStructuresAuthor}.  In the second
case the strategy is ``staying alive''.

\begin{hbcorollary}{thm:p03-w03-p03-064-alice-complete}{proved=annex:p03-w03-p03-064-alice-complete}{The Alice's Restaurant theory is complete}
Let \(T\) be the theory of \cref{def:p03-w03-p03-067-schema}: irreflexivity,
symmetry and all extension axioms \(A_{r,s}\).  Then \(T\) has no finite
models, and Duplicator wins \(\mathrm{EHR}[G,G';t]\) on any two countable
models \(G,G'\) of \(T\) by choosing, for each new vertex picked by Spoiler, a
new vertex of the other graph with the same adjacencies to the vertices
already picked.  Hence \(T\) is complete by
\cref{thm:p03-w03-p03-064-complete-game}, and for each constant
\(0<p<1\) every first-order graph sentence holds in \(G(n,p)\) almost surely
or almost never.
\end{hbcorollary}

Spencer's estimate that each \(A_{r,s}\) fails with probability at most
\(n^{r+s}(1-\epsilon)^{n-r-s}\), \(\epsilon=p^r(1-p)^s\), and the
\(\aleph_0\)-categoricity of \(T\) are
\cref{thm:p03-w03-p03-067-witness,thm:p03-w03-p03-067-categorical}
\cite[pp.~12--14]{Spencer2007LogicRandomStructuresAuthor}.  The direct
completeness proof of \cref{thm:p03-w03-p03-067-complete} plays the same
strategy on models of any cardinality; the corollary above follows the
source's route through countable models.

\section{The almost sure theory for \(p=n^{-\alpha}\)}
\label{sec:p03-w03-p03-064-theory}

This section follows Spencer's Section~5
\cite[pp.~9--12]{Spencer2007LogicRandomStructuresAuthor}.  It writes down a
theory \(T_\alpha\), shows that each of its sentences holds almost surely in
\(G(n,n^{-\alpha})\), and shows by the Ehrenfeucht game that \(T_\alpha\) is
complete.  Together these prove the Shelah--Spencer zero--one law, which
\cref{thm:p03-w03-p03-065-shelah-spencer} states.  Throughout,
\(\alpha\in(0,1)\) is irrational and \(p=n^{-\alpha}\).  Spencer leaves out
three parts of the argument: the large deviation estimate for the number of
extensions, the counting of extensions that are not generic, and the
combinatorial lemmas behind Duplicator's strategy
\cite[pp.~10--12]{Spencer2007LogicRandomStructuresAuthor}.  The annex proves
all three.

\begin{hbdefinition}{def:p03-w03-p03-064-vminuse}{The quantity \(v-e\alpha\) of a vertex set}
For a finite graph \(H\) and \(S\subseteq V(H)\) put
\[
  f(S,H):=|V(H)\setminus S|-\alpha\,e_H(S),
\]
where \(e_H(S)\) is the number of edges of \(H\) with at least one end
outside \(S\).  For a rooted graph \((R,H)\), \(f(R,H)\) is Spencer's
\(v-e\alpha\); for \(R\subseteq S\subsetneq V(H)\), \(f(S,H)\) is the
\(v-e\alpha\) of the nailextension \((S,H)\); for \(R\subsetneq S\subseteq
V(H)\), \(f(R,H|_S)\) is the \(v-e\alpha\) of the subextension
\((R,H|_S)\); and \(f(V(H),H)=0\).  So \((R,H)\) is rigid when
\(f(S,H)<0\) for all \(R\subseteq S\subsetneq V(H)\), and safe when
\(f(R,H|_S)>0\) for all \(R\subsetneq S\subseteq V(H)\)
(\cref{def:p03-w03-p03-065-dense-sparse}).
\end{hbdefinition}

\begin{hblemma}{thm:p03-w03-p03-064-f-laws}{proved=annex:p03-w03-p03-064-f-laws}{Laws for \(v-e\alpha\)}
Let \(H\) be a finite graph and \(R\subseteq V(H)\).
\begin{enumerate}
\item \emph{Additivity.}  If \(R\subseteq S\subseteq V(H)\), then
  \(f(R,H)=f(R,H|_S)+f(S,H)\).
\item \emph{Irrationality.}  If \(S\subsetneq V(H)\), then \(f(S,H)\neq0\).
\item \emph{Safe extensions.}  If \((R,H)\) is safe, \(\kappa\) is the least
  value of \(f(R,H|_S)\) over \(R\subsetneq S\subseteq V(H)\), and
  \(R\subsetneq S\subseteq V(H)\), then \(\kappa>0\) and
  \(f(S,H)\leq f(R,H)-\kappa<f(R,H)\).
\item \emph{Nonsafe extensions contain rigid ones.}  If \((R,H)\) with
  \(R\subsetneq V(H)\) is not safe, there is \(S\) with
  \(R\subsetneq S\subseteq V(H)\) such that \((R,H|_S)\) is rigid.
\end{enumerate}
\end{hblemma}

Item~4 is the combinatorial lemma whose proof Spencer omits in the
completeness argument \cite[p.~12]{Spencer2007LogicRandomStructuresAuthor}.
It also shows that Spencer's two forms of the Nonexistence schema agree:
a graph contains a finite subgraph with \(v-e\alpha<0\) exactly when
\(cl_t(\emptyset)\neq\emptyset\) for some \(t\geq1\)
(\cref{thm:p03-w03-p03-064-closure-models}).

\begin{hbdefinition}{def:p03-w03-p03-064-generic}{\(t\)-generic extensions}
Let \(t\geq1\), let \(G\) be a graph, and let \(x_1,\ldots,x_r\) and
\(y_1,\ldots,y_v\) be distinct vertices.  The tuple \(y_1,\ldots,y_v\) is
\emph{\(t\)-generic} over \(x_1,\ldots,x_r\) if, whenever distinct vertices
\(z_1,\ldots,z_w\) with \(w\leq t\), distinct from all \(x_i\) and \(y_j\),
form a rigid extension over \(x_1,\ldots,x_r,y_1,\ldots,y_v\), there is no
edge between any \(z_k\) and any \(y_j\)
\cite[p.~10]{Spencer2007LogicRandomStructuresAuthor}.  Every tuple counts as
\(0\)-generic.
\end{hbdefinition}

Spencer's gloss is that a generic extension acquires no special properties
relative to its roots: every rigid extension over the \(x\)'s and \(y\)'s is
in fact an extension over the \(x\)'s alone
\cite[p.~10]{Spencer2007LogicRandomStructuresAuthor}.  The condition is
first-order: for fixed \(r,v,t\) there are finitely many rigid rooted graphs
with \(r+v\) roots and at most \(t\) nonroots, and for each one the condition
is a universal statement about \(z_1,\ldots,z_w\).

\begin{hbdefinition}{def:p03-w03-p03-064-theory-alpha}{The almost sure theory \(T_\alpha\)}
The theory \(T_\alpha\) in the language \(\{\sim\}\) consists of the graph
axioms (irreflexivity and symmetry) and two schemas
\cite[p.~10]{Spencer2007LogicRandomStructuresAuthor}.
\begin{itemize}
\item \emph{Nonexistence.}  For each finite graph \(H\) with \(v\) vertices
  and \(e\) edges and \(v-e\alpha<0\): there is no copy of \(H\), that is, no
  injective map \(V(H)\to V(G)\) carrying edges to edges.
\item \emph{Generic Extension.}  For each safe rooted graph \((R,H)\) with
  parameters \(r,v,e\) and each \(t\geq0\): for all distinct
  \(x_1,\ldots,x_r\) there exist \(y_1,\ldots,y_v\) such that
  \begin{enumerate}
  \item \(y_1,\ldots,y_v\) form an \((R,H)\) extension over
    \(x_1,\ldots,x_r\);
  \item there are no edges \(y_iy_j\) or \(y_ix_j\) other than those that
    \(H\) requires;
  \item \(y_1,\ldots,y_v\) is \(t\)-generic over \(x_1,\ldots,x_r\).
  \end{enumerate}
\end{itemize}
\end{hbdefinition}

\begin{hbexample}{def:p03-w03-p03-064-generic-example}{A generic common neighbour at \(\alpha=\pi/7\)}
Take the common-neighbour rooted graph of item~1 of
\cref{def:p03-w03-p03-065-pi-seven}: roots \(x_1,x_2\), one nonroot \(y\)
adjacent to both, \(f=1-2\alpha>0\), safe.  At \(\alpha=\pi/7\),
\(2\alpha<1<3\alpha\), so one vertex \(z\) is rigid over
\(\{x_1,x_2,y\}\) exactly when it is adjacent to all three.  Hence a common
neighbour \(y\) of \(x_1,x_2\) is \(1\)-generic exactly when \(x_1,x_2,y\)
have no common neighbour.  The sentence of \(T_\alpha\) for this rooted graph
and \(t=1\) says that every two distinct vertices have a common neighbour
\(y\) such that \(x_1,x_2,y\) have no common neighbour; condition~2 adds
nothing here, because both possible edges \(yx_1,yx_2\) are required.
\end{hbexample}

For distinct \(\bar x=(x_1,\ldots,x_r)\) let \(N(\bar x)\) be the number of
\((R,H)\) extensions over \(\bar x\), counted as ordered tuples
\(y_1,\ldots,y_v\), and let \(\mu=\mathbb E[N(\bar x)]=(n-r)(n-r-1)\cdots
(n-r-v+1)\,p^e\), which does not depend on \(\bar x\).

\begin{hbtheorem}{thm:p03-w03-p03-064-large-deviation}{proved=annex:p03-w03-p03-064-large-deviation}{Large deviation and counting for safe extensions}
Let \((R,H)\) be safe with parameters \(r,v,e\), and let \(\varepsilon>0\).
\begin{enumerate}
\item For every fixed \(\bar x\),
  \(\Pr\bigl[|N(\bar x)-\mu|>\varepsilon\mu\bigr]=o(n^{-r})\).
\item \emph{Counting Theorem.}  Almost surely
  \((1-\varepsilon)\mu\leq N(\bar x)\leq(1+\varepsilon)\mu\) for all
  distinct \(x_1,\ldots,x_r\).
\end{enumerate}
\end{hbtheorem}

Spencer states both and omits the proof of the first, describing it as the
one technical part of the probability analysis
\cite[pp.~10--11]{Spencer2007LogicRandomStructuresAuthor}.  The difficulty
he names is that the indicators of the potential extensions overlap, so
\(N(\bar x)\) is not binomial.  The annex proof bounds a high central moment
of \(N(\bar x)\) by grouping overlapping extensions.  Spencer adds that the
probability in item~1 can be bounded by an exponential of a power of \(n\)
for a positive exponent \(\lambda\), and that \(o(n^{-r})\) suffices for the
theory \cite[p.~10]{Spencer2007LogicRandomStructuresAuthor}.  The bound is
printed there as \(\exp[n^{-\lambda}]\), which exceeds \(1\); the bound
meant, and proved next, is \(\exp[-n^{\lambda}]\).

\begin{hbcorollary}{thm:p03-w03-p03-064-exp-deviation}{proved=annex:p03-w03-p03-064-exp-deviation}{The deviation probability is at most \(\exp[-n^{\lambda}]\)}
Let \((R,H)\) be safe with parameters \(r,v,e\) and \(v\geq1\), let
\(\kappa\) be as in item~3 of \cref{thm:p03-w03-p03-064-f-laws}, and let
\(\varepsilon>0\).  Put \(\lambda=\kappa/(8v)\).  There is \(n_0\) such that
for all \(n\geq n_0\) and all distinct \(x_1,\ldots,x_r\),
\[
  \Pr\bigl[|N(\bar x)-\mu|>\varepsilon\mu\bigr]\leq\exp[-n^{\lambda}].
\]
Consequently the probability that some \(\bar x\) has
\(|N(\bar x)-\mu|>\varepsilon\mu\) is at most \(n^r\exp[-n^{\lambda}]\) for
\(n\geq n_0\).
\end{hbcorollary}

The proof keeps the moment bound of the annex proof of
\cref{thm:p03-w03-p03-064-large-deviation}, which holds for every order
\(2k\), and lets \(k\) grow like a small power of \(n\).  The exponent
\(\lambda\) is not optimized.

\begin{hblemma}{thm:p03-w03-p03-064-moment-bound}{proved=annex:p03-w03-p03-064-moment-bound}{An upper bound for counts of arbitrary extensions}
Let \((R,J)\) be any rooted graph with \(r\) roots and \(w\) nonroots, and
let \(F^+\) be the larger of \(0\) and the greatest \(f(S,J)\) over
\(R\subseteq S\subsetneq V(J)\).  Let \(N_J(\bar x)\) be the number of
\((R,J)\) extensions over \(\bar x\).  For every \(\delta>0\) and integer
\(k\geq1\) there is \(C\) such that, for all \(n\) and all distinct
\(\bar x\), \(\Pr[N_J(\bar x)\geq n^{F^++\delta}]\leq Cn^{-k\delta}\).
\end{hblemma}

\begin{hbtheorem}{thm:p03-w03-p03-064-theory-as}{proved=annex:p03-w03-p03-064-theory-as}{Every sentence of \(T_\alpha\) holds almost surely}
Each sentence of the Nonexistence schema and of the Generic Extension schema
holds almost surely in \(G(n,n^{-\alpha})\).  Moreover, for a safe
\((R,H)\) and \(t\geq0\), almost surely for every \(\bar x\) all but
\(O(n^{f(R,H)-\delta})\) of the \((R,H)\) extensions over \(\bar x\) satisfy
conditions 2 and 3 of the Generic Extension schema, for some \(\delta>0\)
depending on \((R,H)\) and \(t\).
\end{hbtheorem}

The Nonexistence schema is \cref{thm:p03-w03-p03-064-nonexistence}.  For
Generic Extension, Spencer's outline is that an extension which is not
\(t\)-generic comes with a rigid extension, which lowers \(v-e\alpha\), so
such extensions are counted by a smaller power of \(n\) than \(\mu\)
\cite[p.~11]{Spencer2007LogicRandomStructuresAuthor}.  The annex makes the
lowering exact: for each of the finitely many ways an extension can fail
conditions 2 or 3, the enlarged rooted graph \((R,J)\) has
\(F^+<f(R,H)\), and \cref{thm:p03-w03-p03-064-moment-bound} applies.

\begin{hblemma}{thm:p03-w03-p03-064-closure-models}{proved=annex:p03-w03-p03-064-closure-models}{Closures in models of the Nonexistence schema}
Let \(G\) be a graph with no finite subgraph \(F\) with
\(|V(F)|-|E(F)|\alpha<0\), and let \(t\geq1\).  Write \(\epsilon_t\) for the
\(\epsilon\) of \cref{thm:p03-w03-p03-064-finite-closure}.
\begin{enumerate}
\item \(cl_t(\emptyset)=\emptyset\), and for every finite \(X\subseteq V(G)\),
  \(cl_t(X)\) is finite, is reached from \(X\) by finitely many steps, and
  \(|cl_t(X)\setminus X|<|X|/\epsilon_t\).
\item Conversely, a graph that contains a finite subgraph \(F\) with
  \(|V(F)|-|E(F)|\alpha<0\) has \(cl_t(\emptyset)\neq\emptyset\) for
  \(t=|V(F)|\).
\end{enumerate}
In any graph \(G\), for \(t\geq0\), if \(X\subseteq W\subseteq V(G)\), \(X\)
is finite and \(W\) is \(t\)-closed, then the \(t\)-closure of \(X\) in \(G\)
equals its \(t\)-closure in \(G|_W\).
\end{hblemma}

Duplicator's strategy uses a sequence of closure parameters chosen in advance.
Fix the number \(u\geq1\) of rounds.  Put \(t_0=0\), \(t_1=1\) and, for
\(2\leq i\leq u\),
\[
  t_i=\max\bigl(t_{i-1},\ \lceil u/\epsilon_{t_{i-1}}\rceil+1\bigr).
\]
Spencer requires \(t_{i+1}\geq t_i\) and that \(t_{i+1}\) bound the number of
vertices that a \(t_i\)-closure adds, which exists by the Finite Closure
Lemma \cite[p.~11]{Spencer2007LogicRandomStructuresAuthor}.  The choice above
gives that bound in every model of \(T_\alpha\) by
\cref{thm:p03-w03-p03-064-closure-models}, for all sets of at most \(u\)
vertices.

\begin{hbtheorem}{thm:p03-w03-p03-064-duplicator}{proved=annex:p03-w03-p03-064-duplicator}{Duplicator wins on models of \(T_\alpha\)}
Let \(G,G'\) be models of \(T_\alpha\) and \(u\geq1\).  Duplicator wins
\(\mathrm{EHR}[G,G';u]\).  A winning strategy: after round \(s\), with
\(i=u-s\) rounds remaining, she keeps an isomorphism from
\(cl_{t_i}(x_1,\ldots,x_s)\) in \(G\) onto \(cl_{t_i}(x_1',\ldots,x_s')\) in
\(G'\) sending each \(x_j\) to \(x_j'\).  Write \(\mathit{BIG}=t_i\) and
\(\mathit{SMALL}=t_{i-1}\) before round \(s+1\), and suppose Spoiler selects
\(y\) in \(G\) (the case of \(G'\) is symmetric).
\begin{itemize}
\item \emph{Inside}: if \(y\in cl_{\mathit{BIG}}(x_1,\ldots,x_s)\), she
  selects the image of \(y\) under the isomorphism.
\item \emph{Outside}: otherwise let
  \(\mathit{OLD}=cl_{\mathit{BIG}}(x_1,\ldots,x_s)\) and \(\mathit{NEW}\) the
  set of vertices of \(cl_{\mathit{SMALL}}(x_1,\ldots,x_s,y)\) outside
  \(\mathit{OLD}\).  The rooted graph \((\mathit{OLD},G|_{\mathit{OLD}\cup
  \mathit{NEW}})\) is safe.  She uses the Generic Extension sentence for it,
  with \(t=\mathit{SMALL}\), at the image of \(\mathit{OLD}\), to find
  \(\mathit{NEW}'\) with exactly the same adjacencies, and selects the vertex
  of \(\mathit{NEW}'\) corresponding to \(y\).
\end{itemize}
\end{hbtheorem}

This is Spencer's strategy \cite[pp.~11--12]{Spencer2007LogicRandomStructuresAuthor}.
At the start, \(cl_{t_u}(\emptyset)=\emptyset\) in both graphs by the
Nonexistence schema; at the end, an isomorphism of the \(0\)-closures sending
each \(x_j\) to \(x_j'\) is exactly the winning condition.  The annex proves
the step, including the two facts Spencer leaves to omitted lemmas: that
\(\mathit{NEW}\) over \(\mathit{OLD}\) is safe, and that, by genericity, the
\(\mathit{SMALL}\)-closure of the new position in \(G'\) contains nothing
beyond the copy of the one in \(G\).

\begin{hbcorollary}{thm:p03-w03-p03-064-sparse-law}{proved=annex:p03-w03-p03-064-sparse-law}{\(T_\alpha\) is complete; the zero--one law}
\(T_\alpha\) is consistent, has no finite models and is complete.  For every
first-order sentence \(A\) in the language of graphs, \(A\) holds almost
surely in \(G(n,n^{-\alpha})\) if \(T_\alpha\vdash A\), and almost never
otherwise.  In particular \(\lim_n\Pr[G(n,n^{-\alpha})\entails A]\) exists
and is \(0\) or \(1\).
\end{hbcorollary}

\begin{hbwarning}{def:p03-w03-p03-064-sparse-limits}{What the completeness proof does not give}
The parameters \(t_i\) depend on the number \(u\) of rounds, so the strategy
is one for each finite game; Duplicator need not win the infinite game, and
\(T_\alpha\) is not \(\aleph_0\)-categorical
(\cref{thm:p03-w03-p03-067-sparse-models,thm:p03-w03-p03-064-infinite-game}).
The proof uses irrationality of \(\alpha\) in the dense--sparse dichotomy, in
item~4 of \cref{thm:p03-w03-p03-064-f-laws} and in the closure bounds; at
rational \(\alpha\) the classification breaks down
(\cref{def:p03-w03-p03-065-rational-boundary}).  The corollary does not say which of the two
limits a given sentence has; that is the question whether \(T_\alpha\vdash A\).
\end{hbwarning}

\section{Containment thresholds at every rational exponent}
\label{sec:p03-w03-p03-064-rational-thresholds}

Spencer justifies the claim that the theories \(T_\alpha\) are pairwise
distinct by a known fact: between any two exponents lies a rational \(a/b\),
and some graph \(H\) has the property that the existence of a copy of \(H\)
has threshold function \(n^{-a/b}\)
\cite[p.~15]{Spencer2007LogicRandomStructuresAuthor}.  He gives neither the
graph nor the threshold proof.  This section supplies both, for every
rational \(a/b\in(0,1)\).  The distinctness itself is also proved, by an
extension statement, in \cref{thm:p03-w03-p03-068-distinct}.

\begin{hbdefinition}{def:p03-w03-p03-064-max-density}{Copies of a graph and maximum density}
Let \(H\) be a finite graph with \(v\geq1\) vertices and \(e\) edges.  A
\emph{copy} of \(H\) in a graph \(G\) is an injective map \(\varphi\) from
\(V(H)\) to \(V(G)\) with \(\varphi(x)\sim\varphi(y)\) for every edge \(xy\)
of \(H\); further adjacencies are allowed.  The first-order sentence
\(A_H\) says that a copy exists: it has \(v\) existential quantifiers, the
inequalities between the quantified variables and one adjacency for each
edge of \(H\).  The \emph{maximum density} of \(H\) is
\[
  m(H)=\max\Bigl\{\frac{e_K}{v_K}:K\text{ a subgraph of }H,\ v_K\geq1\Bigr\},
\]
where \(v_K\) and \(e_K\) are the numbers of vertices and edges of \(K\).
The maximum exists because \(H\) has finitely many subgraphs, and it is
attained by an induced subgraph.  If \(H\) has an edge, then
\(m(H)\geq\frac12\).
\end{hbdefinition}

\begin{hbtheorem}{thm:p03-w03-p03-064-containment}{proved=annex:p03-w03-p03-064-containment}{Threshold for containing a fixed graph}
Let \(H\) be a finite graph with at least one edge.  Then \(n^{-1/m(H)}\) is
a threshold function for \(A_H\) in \(G(n,p(n))\)
(\cref{def:p03-w03-p03-065-threshold}): if \(p\ll n^{-1/m(H)}\), then
\(A_H\) holds almost never, and if \(p\gg n^{-1/m(H)}\), then \(A_H\) holds
almost surely.
\end{hbtheorem}

The first half is a first-moment bound applied to a densest subgraph.  The
second is a second-moment bound for the number of copies of \(H\), in which
two copies that overlap on \(s\) vertices share at most \(m(H)s\) edges.
For \(H=K_4\) and \(H=K_5\), whose maximum densities are \(3/2\) and \(2\),
the theorem gives the thresholds \(n^{-2/3}\) and \(n^{-1/2}\) behind
Spencer's spectra \(\{2/3\}\) and \(\{1/2\}\)
\cite[pp.~15--16]{Spencer2007LogicRandomStructuresAuthor}; compare
\cref{thm:p03-w03-p03-065-clique}.

\begin{hbdefinition}{def:p03-w03-p03-064-hab}{The graphs \(H_{a,b}\)}
Let \(a<b\) be positive integers.  Write \(b=qa+r\) with integers \(q\geq1\)
and \(0\leq r<a\), let \(M\) be the least positive integer with
\(aM>2q+2\), and put \(v=aM\).  Let
\[
  P=\bigl\{i\in\{0,1,\ldots,v-1\}:\ \lfloor(i+1)r/a\rfloor-\lfloor ir/a\rfloor=1\bigr\}.
\]
The graph \(H_{a,b}\) has vertex set \(\mathbb Z_v=\{0,1,\ldots,v-1\}\) with
arithmetic modulo \(v\).  Its edges are the pairs \(\{i,i+j\}\) for
\(i\in\mathbb Z_v\) and \(1\leq j\leq q\), together with the chords
\(\{i,i+q+1\}\) for \(i\in P\).
\end{hbdefinition}

So \(H_{a,b}\) is the \(q\)-th power of a \(v\)-cycle with one further chord
of length \(q+1\) at each vertex of \(P\), and \(P\) contains exactly \(r\)
vertices in each of the \(M\) consecutive runs of \(a\) vertices.

\begin{hblemma}{thm:p03-w03-p03-064-hab}{proved=annex:p03-w03-p03-064-hab}{The maximum density of \(H_{a,b}\)}
For positive integers \(a<b\), the graph \(H_{a,b}\) has \(aM\) vertices and
\(bM\) edges, every vertex set \(U\) of \(H_{a,b}\) spans at most
\((b/a)|U|\) edges, and \(m(H_{a,b})=b/a\).
\end{hblemma}

The annex proves the bound by splitting each edge into two nonnegative
weights with sum \(1\), one at each end, so that every vertex receives exactly
\(b/a\).  The chords are not spread evenly by themselves; the weights on the
edges \(\{i,i+1\}\) carry the surplus around the cycle, and the choice of
\(P\) keeps those weights inside \([0,1]\).

\begin{hbcorollary}{thm:p03-w03-p03-064-rational-threshold}{proved=annex:p03-w03-p03-064-rational-threshold}{A threshold at every rational exponent}
For every rational \(a/b\in(0,1)\), with \(a<b\) positive integers, the
existence of a copy of \(H_{a,b}\) has threshold function \(n^{-a/b}\).
Consequently, for irrational \(\alpha<\beta\) in \((0,1)\) and any rational
\(a/b\in(\alpha,\beta)\), \(T_\alpha\vdash A_{H_{a,b}}\) and
\(T_\beta\vdash\neg A_{H_{a,b}}\); in particular \(T_\alpha\) and
\(T_\beta\) have different consequences.
\end{hbcorollary}

\begin{hbexample}{def:p03-w03-p03-064-hab-example}{The graph \(H_{2,3}\)}
For \(a/b=2/3\): \(q=1\), \(r=1\), \(M=3\) (the least \(M\) with
\(2M>4\)) and \(v=6\).  Then \(P=\{1,3,5\}\), and \(H_{2,3}\) is the
\(6\)-cycle \(0,1,\ldots,5\) with the chords \(\{1,3\},\{3,5\},\{5,1\}\):
six vertices and nine edges, density \(3/2\).  In the weighting of the annex,
the chord \(\{i,i+2\}\) gives its whole weight to \(i\) (so \(\{5,1\}\) gives
it to \(5\)), and the cycle edge \(\{i,i+1\}\) gives \(w_i\) to \(i+1\) and
\(1-w_i\) to \(i\), with \((w_0,\ldots,w_5)=(\tfrac12,1,\tfrac12,1,\tfrac12,1)\).
Vertex \(1\) receives \(1\) from \(\{1,3\}\), \(w_0=\frac12\) from
\(\{0,1\}\) and \(1-w_1=0\) from \(\{1,2\}\); vertex \(2\) receives \(w_1=1\)
and \(1-w_2=\frac12\).  Every vertex receives \(\frac32\), so no vertex set
spans more than \(\frac32\) edges per vertex: the odd triangle spans \(3\)
edges, and deleting the vertex \(0\) of degree \(2\) leaves \(7\) edges on
five vertices, below \(7.5\).  By
\cref{thm:p03-w03-p03-064-rational-threshold}, copies of \(H_{2,3}\) appear
at \(n^{-2/3}\), the threshold of \(K_4\), whose maximum density is also
\(3/2\).
\end{hbexample}

\begin{hbwarning}{def:p03-w03-p03-064-threshold-scope}{What the rational thresholds do not give}
The threshold statement concerns \(p\ll n^{-a/b}\) and \(p\gg n^{-a/b}\)
only.  At \(p=n^{-a/b}\) itself the expected number of copies of a densest
subgraph is of constant order, and the theorem says nothing there
(\cref{def:p03-w03-p03-065-threshold-limits}).  The graph \(H_{a,b}\) is one
choice among many; its role is to show that every rational in \((0,1)\)
occurs as a containment threshold, which is what Spencer's argument for
distinctness uses.
\end{hbwarning}

\section{Representing every relation in the random ordered graph}
\label{sec:p03-w03-p03-064-k-ary}

In the random ordered graph \(G_<(n,\frac12)\) of
\cref{def:p03-w03-p03-071-gord}, Spencer represents addition and
multiplication on \(\{1,\ldots,u\}\) by witness vertices whose trace on
\(\{1,\ldots,u\}\) is the triple in question.  He remarks that further work,
perhaps with a changed bound on \(u\), would represent every ternary, indeed
every \(k\)-ary, relation almost surely
\cite[p.~23]{Spencer2007LogicRandomStructuresAuthor}.  The corrected
representation of \(+\) and \(\ast\) is
\cref{thm:p03-w03-p03-071-representation}.  A trace is a set, so it forgets
the order and the repetitions of a tuple.  The encoding below therefore
spends \(k+1\) consecutive vertices on each tuple: one with empty trace that
marks the start, then one vertex per coordinate whose trace is that
coordinate alone.

\begin{hbdefinition}{def:p03-w03-p03-064-rep}{The block encoding \(\operatorname{Rep}^k_u\)}
Fix \(k\geq1\).  In an ordered graph on \(\Omega=\{1,\ldots,n\}\), with
\(y=x+1\) the successor formula and \(N_u(e)\) the trace of \(e\) on
\(\{1,\ldots,u\}\) (\cref{def:p03-w03-p03-071-gord}), let
\(\operatorname{Rep}^k_u(c,d;x_1,\ldots,x_k)\) be the formula
\begin{align*}
  &\textstyle\bigwedge_{i=1}^k x_i\leq u\ \wedge\ \exists_{e_0}\cdots\exists_{e_k}\Bigl[
    c\leq e_0\wedge e_k<d\\
  &\qquad\textstyle\wedge\bigwedge_{i<k}e_{i+1}=e_i+1
    \wedge\forall_w\bigl(w\leq u\to\neg\,w\sim e_0\bigr)\\
  &\qquad\textstyle\wedge\bigwedge_{i=1}^k\bigl(e_i\sim x_i\wedge
    \forall_w\,(w\leq u\wedge w\sim e_i\to w=x_i)\bigr)\Bigr],
\end{align*}
with free variables \(u,c,d,x_1,\ldots,x_k\); here \(x\leq y\) abbreviates
\(x<y\vee x=y\).  It says that some \(e_0\) in \([c,d)\) has empty trace and
is followed by \(k\) consecutive vertices inside \([c,d)\) whose traces are
\(\{x_1\},\ldots,\{x_k\}\).

For \(R\subseteq\{1,\ldots,u\}^k\) with \(m\geq1\) elements listed as
\(\bar r_0,\ldots,\bar r_{m-1}\), a \emph{block layout of \(R\) at \(c\)} is
the condition that for every \(j<m\),
\(N_u(c+j(k+1))=\emptyset\) and \(N_u(c+j(k+1)+i)=\{r_{j,i}\}\) for
\(1\leq i\leq k\), where \(r_{j,i}\) is the \(i\)-th coordinate of
\(\bar r_j\).
\end{hbdefinition}

\begin{hblemma}{thm:p03-w03-p03-064-rep-layout}{proved=annex:p03-w03-p03-064-rep-layout}{A block layout represents its relation}
Fix an ordered graph on \(\Omega\), integers \(k,u\geq1\) and
\(R\subseteq\{1,\ldots,u\}^k\).
\begin{enumerate}
\item If \(R\) has \(m\geq1\) elements, \(u<c\), \(d=c+m(k+1)\leq n\) and
  there is a block layout of \(R\) at \(c\), then for all
  \(\bar x\in\{1,\ldots,u\}^k\), \(\operatorname{Rep}^k_u(c,d;\bar x)\)
  holds if and only if \(\bar x\in R\).
\item If \(R=\emptyset\) and \(u<n\), then
  \(\operatorname{Rep}^k_u(u+1,u+1;\bar x)\) fails for every \(\bar x\).
\end{enumerate}
\end{hblemma}

\begin{hbtheorem}{thm:p03-w03-p03-064-every-relation}{proved=annex:p03-w03-p03-064-every-relation}{Every \(k\)-ary relation is represented almost surely}
Fix \(k\geq1\) and put
\(U_k(n)=\bigl\lfloor(\log_2n/(2k+2))^{1/(k+1)}\bigr\rfloor\).  In
\(G_<(n,\frac12)\), with probability tending to \(1\) as \(n\to\infty\),
for every \(u\) with \(1\leq u\leq U_k(n)\) and every relation
\(R\subseteq\{1,\ldots,u\}^k\) there are \(c\leq d\) in \(\Omega\) with
\(u<c\) such that, for all \(\bar x\in\{1,\ldots,u\}^k\),
\(\operatorname{Rep}^k_u(c,d;\bar x)\) holds exactly when \(\bar x\in R\).
\end{hbtheorem}

For \(k=3\) this represents all ternary relations on \(\{1,\ldots,u\}\) at
once, among them \(\mathrm{Add}_u\) and \(\mathrm{Mul}_u\), for
\(u\leq(\log_2n/8)^{1/4}\).  The changed bound on \(u\) that Spencer
anticipates is visible here: a relation may have \(u^k\) tuples, each costing
\(k+1\) vertices with \(u\) prescribed adjacencies apiece, so a layout has
probability as small as \(2^{-(k+1)u^{k+1}}\), against \(2^{-u\cdot u^2}\)
for the at most \(u^2\) triples of addition in
\cref{thm:p03-w03-p03-071-representation}.

\begin{hbexample}{def:p03-w03-p03-064-rep-example}{The order on \(\{1,2\}\)}
Take \(k=2\), \(u=2\) and \(R=\{(1,2)\}\), the strict order on
\(\{1,2\}\).  A block layout at \(c\geq3\) asks that \(c\) be adjacent to
neither \(1\) nor \(2\), that \(c+1\) be adjacent to \(1\) and not to \(2\),
and that \(c+2\) be adjacent to \(2\) and not to \(1\): six prescribed
adjacencies, of probability \(2^{-6}\) in \(G_<(n,\frac12)\).  With
\(d=c+3\), the only \(e_0\in[c,d)\) with empty trace is \(c\) (the other two
have nonempty traces), so \(\operatorname{Rep}^2_2(c,c+3;x_1,x_2)\) holds
exactly for \((x_1,x_2)=(1,2)\).  The pair \((2,1)\) would need a vertex with
trace \(\{2\}\) directly after the marker, and \((1,1)\) two consecutive
vertices with trace \(\{1\}\); the layout has neither.
\end{hbexample}

\begin{hbwarning}{def:p03-w03-p03-064-rep-limits}{What the block encoding needs}
The encoding uses the built-in order twice: to find consecutive vertices,
and to separate the witness interval from \(\{1,\ldots,u\}\).  Without an
order, as in \(G(n,n^{-1/4})\), common-neighbourhood relations are symmetric
in their arguments, and the step from symmetric relations to \(<\) and \(+\)
is not available (\cref{def:p03-w03-p03-071-sparse-outline}).  The theorem
gives no upper limit for the \(u\) at which every relation is represented;
for \(+\) and \(\ast\) such a limit is \cref{thm:p03-w03-p03-071-upper}.
\end{hbwarning}

\section{Total influence of bounded-depth circuits}
\label{sec:p03-w03-p03-064-circuits}

Shelah's slow-oscillation theorem for \(G_<(n,\frac12)\)
(\cref{thm:p03-w03-p03-071-shelah}) rests on a bound for the total influence
of bounded-depth circuits.  \Cref{thm:p03-w03-p03-071-ac0-influence} gives
such a bound with \(\log(2Ns)\) in place of \(\log s\), where \(N\) is the
number of inputs, which suffices there because the circuits have size
polynomial in \(N\).  Boppana's theorem, in the form O'Donnell states as his
Theorem~4.30 and attributes to Boppana, has no dependence on \(N\)
\cite[thm.~4.30, p.~106; p.~111]{ODonnell2021AnalysisBooleanArxiv}.  This
section proves that form.  It uses the circuits, restrictions, decision
trees and normal forms of
\cref{def:p03-w03-p03-071-circuit,def:p03-w03-p03-071-restriction}, the
restriction facts \cref{thm:p03-w03-p03-071-restriction-facts} and the
switching lemma \cref{thm:p03-w03-p03-071-switching}.  A depth-\(d\) circuit
has \emph{width} at most \(w\) if every layer-\(1\) gate has at most \(w\)
inputs \cite[def.~4.27, p.~104]{ODonnell2021AnalysisBooleanArxiv}.

\begin{hblemma}{thm:p03-w03-p03-064-influence-basics}{proved=annex:p03-w03-p03-064-influence-basics}{Two facts about total influence}
Let \(f:\{0,1\}^V\to\{0,1\}\) with \(V\) finite.
\begin{enumerate}
\item If \(f\) is the AND or the OR of \(g_1,\ldots,g_m\), and coordinate
  \(b\) is pivotal for \(f\) at \(x\), then \(b\) is pivotal for some \(g_i\)
  at \(x\).  Hence
  \(\operatorname{Inf}_b[f]\leq\sum_i\operatorname{Inf}_b[g_i]\).
\item If \(f\) is computed by a DNF or a CNF of width at most \(w\), then
  \(\mathbf I[f]\leq2w\).
\end{enumerate}
\end{hblemma}

Item~2 is the point at which the number of inputs disappears: a narrow
normal form has small total influence however many variables it reads.

\begin{hbtheorem}{thm:p03-w03-p03-064-depth-induction}{proved=annex:p03-w03-p03-064-depth-induction}{Total influence of narrow bounded-depth circuits}
Put \(c_2=2\) and \(c_d=32(2^{d-2}c_{d-1}+1)\) for \(d\geq3\).  Let
\(d\geq2\), \(w\geq1\) and \(S\geq1\), and let \(f\) be computed by a
depth-\(d\) circuit of width at most \(w\) with at most \(S\) gates in
layers \(2,\ldots,d-1\).  Then
\[
  \mathbf I[f]\leq c_d\,w\,\Lambda^{d-2},\qquad \Lambda=\log_2(4wS).
\]
\end{hbtheorem}

\begin{hbtheorem}{thm:p03-w03-p03-064-boppana}{proved=annex:p03-w03-p03-064-boppana}{Total influence of bounded-depth circuits (Boppana)}
For \(d\geq2\) put \(K_d=3^{d-1}c_{d+1}\).  If \(f\) is computed by a
depth-\(d\) circuit of size \(s\geq2\), then
\(\mathbf I[f]\leq K_d(\log_2s)^{d-1}\).
\end{hbtheorem}

Changing the base of the logarithm changes only \(K_d\).  The proof treats
the circuit as one of depth \(d+1\) and width \(1\), by placing a one-input
gate on each literal, so that the unbounded fan-in of layer \(1\) is
absorbed by the first random restriction; the exponent \(d-1\) is
\((d+1)-2\).  At each depth, the event that the switching lemma fails is paid
for with item~1 of \cref{thm:p03-w03-p03-064-influence-basics}, which bounds
the influence of the whole circuit by that of its layer-\(2\) gates, instead
of with the number of free variables.  The constants are far from optimal.

\begin{hbexample}{def:p03-w03-p03-064-boppana-example}{The two bounds on one circuit}
Let \(f\) be the OR of the \(N\) variables, a depth-\(2\) circuit (an OR of
\(N\) one-literal AND gates) of size \(s=N\) and width \(1\).  Coordinate
\(b\) is pivotal exactly at the input in which all other coordinates are
\(0\), so \(\operatorname{Inf}_b[f]=2^{1-N}\) and
\(\mathbf I[f]=N2^{1-N}\).  \Cref{thm:p03-w03-p03-064-influence-basics},
item~2, gives \(\mathbf I[f]\leq2\), independent of \(N\).
\Cref{thm:p03-w03-p03-064-boppana} gives \(K_2\log_2N\), and
\cref{thm:p03-w03-p03-071-ac0-influence} a bound of order \(\log(2N^2)\).
Majority on \(N\) inputs has total influence of order \(\sqrt N\)
\cite[p.~107]{ODonnell2021AnalysisBooleanArxiv}, so for fixed \(d\) no
depth-\(d\) circuit of size \(2^{o(N^{1/(2d-2)})}\) computes it, by
\cref{thm:p03-w03-p03-064-boppana}.
\end{hbexample}

\begin{hbwarning}{def:p03-w03-p03-064-circuit-limits}{What the influence bound does not say}
The bound is useful only when \((\log s)^{d-1}\) is small against the number
of inputs.  For \(s\) polynomial in \(N\) it is a power of \(\log N\), which
is what \cref{thm:p03-w03-p03-071-slice} needs to make adjacent slices
indistinguishable.  For depth growing with \(N\), or size exponential in a
power of \(N\), it gives nothing.  It bounds the average sensitivity, which
is the quantity the slice argument uses, and does not by itself decide what a
circuit computes on a particular input.
\end{hbwarning}

\section{Relationships}
\label{sec:p03-w03-p03-064-relationships}

The closure of \cref{sec:p03-w03-p03-064-closure} is built from the rigid
rooted graphs of \cref{def:p03-w03-p03-065-dense-sparse}; the Finite Closure
Lemma is the reason a Duplicator strategy for \(G(n,n^{-\alpha})\) needs to
look only at boundedly many vertices around the chosen ones
\cite[p.~11]{Spencer2007LogicRandomStructuresAuthor}.  The safe-extension
criterion \cref{thm:p03-w03-p03-065-safe-criterion} supplies the existence
half of such a strategy.  The game corollaries of
\cref{sec:p03-w03-p03-064-game} rest on \cref{thm:p03-ef}, and through
\cref{thm:p03-w03-p03-067-as-theory} they turn a zero--one law into a
statement about countable models.  The uniform case \(p=1/2\) and the
uniform relational structures are in \cref{sec:p03-zero-one-law}.

\section{Sources and further reading}
\label{sec:p03-w03-p03-064-sources}

Spencer's closing section, \emph{About the references}
\cite[pp.~27--28]{Spencer2007LogicRandomStructuresAuthor}, points the reader
to the following literature; the bibliographic data below are those of his
reference list \cite[pp.~28--30]{Spencer2007LogicRandomStructuresAuthor}.
\begin{itemize}
\item \emph{Surveys}: Compton \cite{Compton1988ZeroOneLaws}, Winkler
  \cite{Winkler1993RandomStructures}, Lynch \cite{Lynch1997RandomFiniteModels}
  and Spencer \cite{Spencer1993VariableProbability}.
\item \emph{The game}: Ehrenfeucht first gave it \cite{Ehrenfeucht1961}; it
  is essentially present in earlier work of Fra\"iss\'e, and is also called
  the Ehrenfeucht--Fra\"iss\'e game.
\item \emph{The uniform case \(p=1/2\)}: the zero--one law is due to
  Glebski\u\i, Kogan, Liogon'ki\u\i\ and Talanov \cite{GlebskiiEtAl1969}
  and to Fagin \cite{Fagin1976}.
\item \emph{Random graphs}: the subject begins with Erd\H{o}s and R\'enyi
  \cite{ErdosRenyi1960Evolution}; the basic text is Bollob\'as
  \cite{Bollobas1985RandomGraphs}.
\item \emph{The sparse law}: the zero--one law for \(p=n^{-\alpha}\) first
  appeared in Shelah and Spencer \cite{ShelahSpencer1988Sparse}; a proof by
  the Ehrenfeucht game is in Spencer \cite{Spencer1991EhrenfeuchtGame}, a
  syntactic proof of completeness of \(T_\alpha\) in Spencer
  \cite{Spencer1991Continuum}, an examination of the countable models of
  \(T_\alpha\) in Spencer \cite{Spencer1990CountableSparse}, and some of this
  material also in Alon and Spencer \cite{AlonSpencer1991ProbabilisticMethod}.
\item \emph{Other random structures}: Lynch on unary functions
  \cite{Lynch1985UnaryFunctions}; Shelah and Spencer
  \cite{ShelahSpencer1994UnaryPredicates} and St.~John and Spencer
  \cite{SpencerStJohn1998UnaryPredicates} on random unary predicates with an
  order, which Spencer stresses differ considerably from his Section~1;
  {\L}uczak on random partial orders \cite{Luczak1991Posets}; and {\L}uczak
  and Shelah on a model on \(\{1,\ldots,n\}\) whose adjacency probability for
  \(i,j\) depends on \(|i-j|\) \cite{LuczakShelah1995Homogeneous}.
\item \emph{Stronger logics}: these mostly yield negative results.  Kaufmann
  and Shelah give a nonconvergent second-order sentence for \(G(n,1/2)\)
  \cite{KaufmannShelah1983Monadic}; many such results, including those on the
  random ordered graph, are in Compton, Henson and Shelah
  \cite{ComptonHensonShelah1987}; Shelah shows that on the random ordered
  graph no first-order sentence almost surely encodes evenness of the size
  \cite{Shelah1995VeryWeak}; and Hella, Kolaitis and Luosto study almost sure
  equivalence of logics in general \cite{HellaKolaitisLuosto1996}.
\item \emph{Extension statements and the range of the law}: Spencer studies
  extension statements in detail \cite{Spencer1990ThresholdExtension};
  {\L}uczak and Spencer nearly characterize the functions \(p=p(n)\), not only
  those of the form \(n^{-\alpha}\), for which the zero--one law holds
  \cite{LuczakSpencer1991WhenZeroOne}; and Spencer and Tardos give necessary
  conditions on the limiting probability of a fixed sentence \(A\) in
  \(G(n,n^{-\alpha})\), viewed as a function of \(\alpha\)
  \cite{SpencerTardosUpsDownsSubmitted},
  where Spencer describes the sufficiency proof by Tardos as in preparation
  \cite[p.~28]{Spencer2007LogicRandomStructuresAuthor}.
\end{itemize}
The bibliographic entries for these works are taken from Spencer's list and
have not been checked against the publications themselves; no page of them
is cited here.  The motto of the survey is a line attributed to Michael Rabin
\cite[p.~1]{Spencer2007LogicRandomStructuresAuthor}.
